%% ==============================================================================
%% backmatter/appendix-a.tex
%% Annexe A : Aide-Mémoire & Guide Pratique des Commandes (Cheat Sheet)
%% ==============================================================================
\chapter{Aide-Mémoire et Guide des Commandes}
\label{anx:aide-memoire}

\section{Récapitulatif des Commandes Métier et Macros}
\label{sec:recap-macros}

Les \cref{tab:recap-math,tab:recap-env-outils} rassemblent les raccourcis sémantiques,
les boîtes visuelles et les commandes d'orchestration définis dans \code{config/macros.sty}.

\begin{table}[htbp]
  \centering
  \caption{Aide-mémoire des notations mathématiques et balisage sémantique textuel}
  \label{tab:recap-math}
  \begin{tabularx}{\textwidth}{@{} l l >{\raggedright\arraybackslash}X @{}}
    \toprule
    \textbf{Commande} & \textbf{Rendu visuel} & \textbf{Usage et signification} \\
    \midrule
    \code{\textbackslash vect\{x\}} & $\vect{x}$ & Vecteur en gras droit/italique \\
    \code{\textbackslash mat\{M\}} & $\mat{M}$ & Matrice formelle en gras droit \\
    \code{\textbackslash setR}, \code{\textbackslash setC} & $\setR$, $\setC$ & Ensembles des réels et complexes \\
    \code{\textbackslash norm\{x\}} & $\norm{x}$ & Norme mathématique délimitée \\
    \code{\textbackslash scal\{u\}\{v\}} & $\scal{u}{v}$ & Produit scalaire délimité \\
    \code{\textbackslash diff} & $\diff x$ & Symbole différentiel droit pour intégrales \\
    \code{\textbackslash term\{mot\}} & \term{mot} & Terme technique introduit pour la première fois \\
    \code{\textbackslash code\{snippet\}} & \code{snippet} & Fragment de code ou nom de fichier en ligne \\
    \bottomrule
  \end{tabularx}
\end{table}

\begin{table}[htbp]
  \centering
  \caption{Aide-mémoire des boîtes graphiques, interactivité et orchestration}
  \label{tab:recap-env-outils}
  \begin{tabularx}{\textwidth}{@{} l l >{\raggedright\arraybackslash}X @{}}
    \toprule
    \textbf{Environnement / Commande} & \textbf{Type} & \textbf{Usage et comportement} \\
    \midrule
    \code{callout} & Boîte cadrée & Remarque, conseil ou note méthodologique \\
    \code{hypothesis} & Boîte numérotée & Hypothèse scientifique (compatible \code{cleveref}) \\
    \code{definitionbox} & Boîte d'en-tête & Définition formelle avec bandeau coloré \\
    \code{boiteSynthese} & Boîte synthèse & Synthèse didactique avec cartouche coloré \\
    \code{cadreFigure}   & Cadre arrondi & Encadrement soigné pour figures et schémas \\
    \code{\textbackslash ocgbutton\{id\}\{txt\}} & Bouton cliquable & Bascule interactive d'un calque vectoriel OCG \\
    \code{\textbackslash ocgstack\{vue1\}\{vue2\}} & Double vue & Superposition écran ou linéarisation sur papier \\
    \code{\textbackslash includepdfpaper\{doc\}} & Pages PDF & Insertion multipage avec en-têtes préservés \\
    \code{\textbackslash geometrieLarge} & Marges 1,5 cm & Pleine page pour grands tableaux ou couverture \\
    \code{\textbackslash geometrieWord}  & Marges 2,5 cm & Profil de révision avec marge latérale pour notes \\
    \code{\textbackslash geometrieStandard} & DIV=12 & Restauration de la géométrie canonique KOMA \\
    \code{\textbackslash afficherLogo}, \code{\textbackslash masquerLogo} & Logo scrlayer & Contrôle du calque flottant institutionnel \\
    \code{\textbackslash todoH[st][col]\{n\}\{txt\}} & Note marge & Troncature, infobulle et lien vers la table \\
    \code{\textbackslash missingfigure[opts]\{txt\}} & Réservation & Espace réservé sans bloquer la chaîne \\
    \code{\textbackslash thesistodolist} & Sommaire & Tableau de bord dynamique des tâches de révision \\
    \bottomrule
  \end{tabularx}
\end{table}

\clearpage

\section{Guide des Clés Typographiques KOMA-Script}
\label{sec:koma-font-keys}

Pour adapter l'esthétique générale de votre ouvrage sans toucher au code des chapitres,
vous pouvez modifier les polices attribuées aux éléments KOMA-Script dans \code{mythesis.cls}
via la macro \code{\textbackslash setkomafont\{élément\}\{spécification\}} :

\begin{table}[htbp]
  \centering
  \caption{Principaux éléments typographiques configurables dans \code{mythesis.cls}}
  \label{tab:koma-elements}
  \begin{tabularx}{\textwidth}{@{} l l >{\raggedright\arraybackslash}X @{}}
    \toprule
    \textbf{Élément KOMA} & \textbf{Réglage par défaut} & \textbf{Partie du document affectée} \\
    \midrule
    \code{disposition}      & \code{\textbackslash bfseries\textbackslash sffamily\textbackslash color\{ThesisPrimary\}} & Tous les titres (chapitres, sections) \\
    \code{chapter}          & \code{\textbackslash Huge\textbackslash bfseries\textbackslash sffamily}                   & Titre principal de chapitre \\
    \code{section}          & \code{\textbackslash Large\textbackslash bfseries\textbackslash sffamily}                  & Titre de section de niveau 1 \\
    \code{subsection}       & \code{\textbackslash large\textbackslash bfseries\textbackslash sffamily\textbackslash color\{ThesisSecondary\}} & Titre de sous-section \\
    \code{pageheadfoot}     & \code{\textbackslash small\textbackslash sffamily\textbackslash color\{ThesisMuted\}}     & Texte des en-têtes et pieds de page \\
    \code{pagination}       & \code{\textbackslash small\textbackslash sffamily\textbackslash bfseries}                 & Numéro de page (folio) \\
    \code{captionlabel}     & \code{\textbackslash bfseries\textbackslash sffamily\textbackslash color\{ThesisPrimary\}} & Libellé des légendes (« Figure 1 ») \\
    \code{descriptionlabel} & \code{\textbackslash bfseries\textbackslash sffamily\textbackslash color\{ThesisPrimary\}} & Entrées des listes de description \\
    \bottomrule
  \end{tabularx}
\end{table}

\clearpage

\section{Commandes Terminal Utiles pour la Compilation}
\label{sec:commandes-terminal}

Depuis la racine du dossier \code{mon-projet/} :
\begin{description}
  \item[\code{latexmk}] Compile l'intégralité du document (passes multiples de Lua\LaTeX{} et Biber).
  Le PDF final est placé dans \code{build/main.pdf}.
  \item[\code{latexmk -pvc}] Mode interactif en direct (*continuous preview*) :
  recompile automatiquement le document à chaque sauvegarde de fichier.
  \item[\code{latexmk -c}] Supprime les fichiers auxiliaires de compilation temporaires.
  \item[\code{latexmk -C}] Supprime la totalité des fichiers générés, y compris le PDF produit dans \code{build/}.
\end{description}
